% =============================================================================
% PROTOCOL BUFFER SCHEMAS
% =============================================================================

\newpage
\section{Protocol Buffer Schema Architecture}

This section documents the Protocol Buffer schema architecture for RPi5$\leftrightarrow$RX72N communication. All schemas follow the \textbf{Boston Dynamics Protobuf Style Guide
} for
  consistency, maintainability,
      and industry best practices.

\subsection{Architecture Overview
 }

\textbf{Overview:} Five gRPC services define the complete API surface for robot control, organized by functional domain. Services communicate via the nanopb-serialized frames described in Section 1.

\noindent\makebox[\textwidth]{
  %
\begin{tikzpicture}[font =\small, node distance=0.4cm,
                     service/.style={draw, thick, rectangle,
                                        minimum width = 6cm,
                                        minimum height = 1.1cm, align = center},
                     common/.style={draw, thick, rectangle,
                                       minimum width = 14cm,
                                       minimum height = 1.1cm, align = center,
                                       fill=cyan!20},
  ] % Left column - Core Services
    \node[service, fill=blue!20](motor) at(-4, 0){\textbf{MotorControlService}\\Velocity commands, encoder streaming};
  \node[service, fill=green!20, below=of motor](telemetry){\textbf{TelemetryService}\\Sensor data, system status};

  % Right column - Support Services
    \node[service, fill=orange!20](battery) at(4, 0){\textbf{BatteryManagementService}\\BMS state, protection, balancing};
  \node[service, fill=yellow!20, below=of battery](config){\textbf{ConfigurationService}\\PID gains, calibration, NVS};

  % Common proto spans bottom
    \node[common, below=0.8cm of telemetry, xshift=4cm](common){\textbf{common.proto}\\Headers, status codes, geometry types};

  % Labels
    \node[above=0.3cm of motor, font=\footnotesize\bfseries]{Real-Time Control};
  \node[above=0.3cm of battery, font=\footnotesize\bfseries]{Device Management};

  % Arrows to common - routed around blocks
    \draw[->, thick, gray](motor.west)-- ++(-0.5, 0) |-(common.west);
  \draw[->, thick, gray](telemetry.south)--(telemetry.south |-common.north);
  \draw[->, thick, gray](battery.east)-- ++(0.5, 0) |-(common.east);
  \draw[->, thick, gray](config.south)--(config.south |-common.north);
  \end{tikzpicture} %
}

\vspace{0.3cm}

\layerbox{keyinsightcolor} { Key Design Principles}
{
  %
\begin{itemize}[leftmargin = 1.5em, topsep = 2pt, itemsep = 1pt]
    \item \textbf{Versioned Package :} All services use \texttt {
    star.v1
 } namespace for API evolution
    \item \textbf{Request/Response Headers:} Every RPC includes tracing, timestamps, and status
    \item \textbf{Streaming Support:} Continuous data uses server streaming (up to 100Hz)
    \item \textbf{Static Allocation:} All messages sized for nanopb constraints (no dynamic memory)
\end{itemize}
}

\textbf{Note:} Naming conventions, unit suffixes, and coding standards are consolidated in Section~\ref{sec:protobuf-style} (Style Guide and Conventions).

\subsection{Service Specifications}
\label{sec:service-specifications}

\subsubsection{MotorControlService}

\textbf{Purpose:} Low-latency differential drive control with real-time encoder feedback.

\textbf{Target Latency:} $<$10ms round-trip for velocity commands.

\begin{table}[H]
\centering
\caption{MotorControlService RPCs}
\label{tab:motor-rpcs}
\begin{tabular} { | l | l | p{6cm} |}
\hline
\textbf{RPC} & \textbf{Type} & \textbf{Description}
\\
\hline SetVelocity &Unary &Set left /
    right wheel velocities(m / s) \\
EmergencyStop &Unary &Immediate stop,
    priority over all commands \\
SetMotorPower &Unary &Direct PWM
        control(bypass PID, testing only) \\
StreamEncoders &Server
        stream &Encoder ticks and velocity at 1 -
        100Hz \\
ControlStream &Bidirectional &Velocity in,
    encoder feedback out \\
\hline
\end{tabular}
\end{table}

\textbf{Key Messages:}
\begin{itemize}[leftmargin = 1.5em, topsep = 4pt, itemsep = 2pt]
    \item \texttt {
  VelocityCommand
}
--Left / right velocity(m / s), sequence number, timestamp
    \item \texttt {
  EncoderData
}
--Motor ID, tick count(341.2 PPR), velocity, timestamp
    \item \texttt {
  MotorStatus
}
--Duty cycle, velocity, temperature, current, fault flags
    \item \texttt {
  PidConfig
}
--Kp, Ki, Kd gains with output saturation limits
\end{itemize}

\subsubsection{TelemetryService}

\textbf{Purpose :} Real -
    time sensor data aggregation and system health monitoring.

\begin{table}[H]
\centering
\caption{TelemetryService RPCs}
\label{tab:telemetry-rpcs}
\begin{tabular} { | l | l | p{6cm} |}
\hline
\textbf{RPC} & \textbf{Type} & \textbf{Description}
\\
\hline GetTelemetry &Unary &Snapshot of all sensor
        data \\
StreamTelemetry &Server stream &Continuous sensor updates at 1 -
    100Hz \\
GetSystemStatus &Unary &Firmware version,
    uptime, connection states \\
\hline
\end{tabular}
\end{table}

\textbf{Key Messages:}
\begin{itemize}[leftmargin = 1.5em, topsep = 4pt, itemsep = 2pt]
    \item \texttt {
  TelemetryData
}
--IMU, GPS, battery \%, WiFi signal, CPU usage, temperature
    \item \texttt {
  ImuData
}
--Pitch / roll / yaw(rad), acceleration(m / s$ ^ 2 $), angular velocity(rad / s)
    \item \texttt {
  GpsData
}
--Lat / lon(degrees), altitude(m), accuracy, satellite count, fix type
    \item \texttt {
  SystemStatus
}
--Connection states, operating mode, free heap, uptime
\end{itemize}

\textbf{Enumerations:
}
\begin{itemize}[leftmargin = 1.5em, topsep = 4pt, itemsep = 2pt]
    \item \texttt {
  RobotMode
}
--IDLE, MANUAL, AUTONOMOUS, MAPPING, EMERGENCY\_STOP
    \item \texttt {
  ConnectionStatus
}
--CONNECTED, DISCONNECTED, CONNECTING, ERROR
    \item \texttt { GpsFix}
--NONE, 2D, 3D, DGPS, RTK\_FLOAT, RTK\_FIXED
\end{itemize}

\subsubsection{BatteryManagementService}

\textbf{Purpose:} BQ7850 BMS monitoring for lithium-ion battery pack safety and management.

\textbf{Hardware:} TI BQ7850 16-cell BMS with I2C interface.

\begin{table}[H]
\centering
\caption{BatteryManagementService RPCs}
\label{tab:battery-rpcs}
\begin{tabular} { | l | l | p{6cm} |}
\hline
\textbf{RPC} & \textbf{Type} & \textbf{Description}
\\
\hline GetBatteryState &Unary &Complete battery snapshot(
    cells, temps, SOC) \\
StreamBatteryState &Server stream &Continuous battery
        monitoring \\
Get
        / SetProtectionThresholds &Unary &OV / UV / OC
        / OT protection limits \\
Enable / DisableCellBalancing &Unary &Per
    - cell passive balancing control \\
ControlFets &Unary &Charge
          / discharge FET enable / disable \\
GetDeviceInfo &Unary &Manufacturer,
    serial, chemistry, versions \\
\hline
\end{tabular}
\end{table}

\textbf{Key Messages:}
\begin{itemize}[leftmargin = 1.5em, topsep = 4pt, itemsep = 2pt]
    \item \texttt {
  BatteryState
}
--Composite of all battery data
    \item \texttt { CellData}
--Individual cell voltages(mV), pack voltage, valid cell count
    \item \texttt {
  BatteryTemperatureData
}
--Thermistor readings(0.1\textdegree C), average
    \item \texttt {
  StateOfChargeData
}
--Relative / absolute SOC(\%), capacity(mAh), cycle count
    \item \texttt {
  ProtectionThresholds
}
--OV / UV voltage limits, OC current limits, OT temperatures
\end{itemize}

\subsubsection{ConfigurationService}

\textbf{Purpose :
} Runtime parameter management with NVS persistence across power cycles.

\begin{table}[H]
\centering
\caption{ConfigurationService RPCs}
\label{tab:config-rpcs}
\begin{tabular} { | l | l | p{6cm} |}
\hline
\textbf{RPC} & \textbf{Type} & \textbf{Description}
\\
\hline GetConfiguration &Unary &
        Read current system configuration \\
SetConfiguration &Unary &Apply
        configuration(optional NVS persist) \\
ResetToDefaults &Unary &Restore
    factory defaults \\
ValidateConfiguration &Unary &Validate without
    applying \\
SaveConfiguration &Unary &Persist in
    - memory config to NVS \\
Get / SetMotorPidConfig &Unary &Per
    - motor PID tuning \\
\hline
\end{tabular}
\end{table}

\textbf{Key Messages:}
\begin{itemize}[leftmargin = 1.5em, topsep = 4pt, itemsep = 2pt]
    \item \texttt {
  SystemConfiguration
}
--Complete config with version and CRC32
    \item \texttt { MotorPidConfig}
--Per - motor PID gains(4 motors supported)
    \item \texttt {
  CurrentSensorCalibration
} -- Offset and scale for current sensing
    \item \texttt{SafetyThresholds}
--Overcurrent, thermal shutdown, cell imbalance limits
    \item \texttt {
  TimingConfig
}
--Motor period, telemetry period, communication timeout
\end{itemize}

\subsection{common.proto-- Shared Types}

All services import \texttt {
  common.proto
} for
  standardized request / response handling.

\begin{table}[H]
\centering
\caption{Standard Headers}
\label{tab:headers}
\begin{tabular} { | l | p{9cm} |}
\hline
\textbf{Message} & \textbf{Fields}
\\
\hline RequestHeader &request\_id(UUID), client\_timestamp, client\_version,
    timeout \\
ResponseHeader &request\_id(echo), server\_timestamp, status,
    error \_message, latency\_us \\
\hline
\end{tabular}
\end{table}

\textbf{Status Codes :} UNKNOWN, OK, INVALID\_REQUEST, INTERNAL\_ERROR,
    NOT\_FOUND, TIMEOUT, INVALID\_STATE,
    ESTOP\_ACTIVE

\textbf{Geometric Types:
}
\begin{itemize}[leftmargin = 1.5em, topsep = 4pt, itemsep = 2pt]
    \item \texttt {
  Vector3
}
--Generic 3D vector(x, y, z)
    \item \texttt { Quaternion}
--Orientation(w, x, y, z)
    \item \texttt { SE2Pose}
--2D pose(x\_m, y\_m, angle\_rad)
    \item \texttt { SE2Velocity}
--2D velocity(vx\_mps, vy\_mps, omega\_rad\_per\_s)
\end{itemize}

\subsection{Code Generation and Integration}

\begin{table}[H]
\centering
\caption{Code Generation Targets}
\label{tab:codegen}
\begin{tabular}{|l|l|l|l|}
\hline
\textbf{Target} & \textbf{Generator} & \textbf{Output} & \textbf{Usage} \\
\hline
Go & protoc-gen-go + grpc & \texttt{gen/go/} & star-gateway(RPi5) \\
\hline
nanopb & nanopb\_generator & \texttt{gen/nanopb/} & RX72N firmware \\
\hline
\end{tabular}
\end{table}

\textbf{nanopb Constraints:} RX72N requires static allocation. Field sizes are configured in \texttt{.options} files:
\begin{itemize}[leftmargin=1.5em,topsep=4pt,itemsep=2pt]
    \item String fields: \texttt{max\_size:64} (request\_id), \texttt{max\_size:128} (error\_message)
    \item Repeated fields: \texttt{max\_count:16} (cell voltages), \texttt{max\_count:4} (motors)
    \item Bytes fields: \texttt{max\_size:1024} (firmware chunks)
\end{itemize}

\textbf{Build Integration:}
\begin{itemize}[leftmargin = 1.5em, topsep = 4pt, itemsep = 2pt]
    \item \textbf{Linting:
}
\texttt{buf lint} enforces style guide compliance
    \item \textbf{Breaking Changes:
}
\texttt{buf breaking} detects API incompatibilities
    \item \textbf{CI Pipeline :} Automatic generation and testing on PR
\end{itemize}
